\section{Exact geometric modes and the Jacobi plateau}
\label{sec:rank}

The \(D=50\) lattice calculation lies in the atom-free regime \(D<M\).
To test Eq.~\eqref{eq:atom-plateau}, we use an independent exact
root-response construction that reaches both sides of the capacity boundary.
For \(N=3\) particles in \(n\) orbitals, the cyclic Laughlin zero-mode
dimension and one-tangent descendant capacity are
\begin{equation}
 D(n)=\frac{n}{n-3}\binom{n-3}{3},\qquad
 M(n)=2n(n-3).
 \label{eq:root-dimensions}
\end{equation}
Two independently sampled one-body response maps are assembled in the
monomial root basis and whitened by a thin QR factorization.
This directly produces the row isometry in
Eq.~\eqref{eq:isometry-compression} as an independent construction from the
physical Kapit--Mueller calculation.

The registered sequence contains \(\RootMatricesTotal\) matrices at
\[
 D=16,\ 50,\ 112,\ 210,\ 352,\ 546,\ 800 .
\]
It is a finite-rank sequence at fixed particle number \(N=3\), designed to
resolve the exact capacity transition.
At \(D=546\), \(M=540\) and the atom theorem requires six levels at each
boundary.
At the endpoint \((D,M)=(\LargestRankVTwo,\LargestChannels)\), it requires
\(\LargestAtomEachVTwo\) levels at each boundary and leaves
\(k=560\) continuously distributed levels.

\begin{figure*}[t]
 \includegraphics[width=\textwidth]{figure_5_jacobi_atoms_v2.pdf}
 \caption{
 \textbf{Exact boundary atoms reshape the connected plateau while the
 continuous sector retains the Jacobi ramp.}
 (a) Full-spectrum connected curvature SFF at the first capacity crossing,
 \(D=546,M=540\), with six exact levels at each of \(\lambda=\pm1\).
 (b) Endpoint \(D=800,M=680\), with \(120\) atoms per boundary.
 The full physical root-response curve follows the full Jacobi prediction,
 while removing the algebraically identified atoms restores the
 unit continuous-sector plateau.
 (c) Across the seven-rank sequence, the exact full connected plateau is
 \(k/D\); it falls as the deterministic atom weight
 \(2(D-M)_+/D\) grows.
 (d) Raw, non-unfolded endpoint SFF separated into atom--atom,
 atom--continuum, and continuum--continuum terms according to
 Eq.~\eqref{eq:atom-raw-decomposition}.
 Their sum closes the full raw form factor to \(1.14\times10^{-13}\).
 }
 \label{fig:jacobi-atoms}
\end{figure*}

The endpoint plateau is therefore fixed analytically:
\begin{equation}
 K_{F,c}^{\rm full}(\infty)
 =\frac{2M-D}{D}
 =\LargestConnectedPlateau .
 \label{eq:endpoint-plateau}
\end{equation}
The reference curve is obtained in two independent ways: evaluation of the
determinantal kernel in Eq.~\eqref{eq:finite-jacobi-sff} and matrix-beta
sampling of the Jacobi ensemble.
The response spectra are generated independently from both references.

This result refines the meaning of random-matrix-like curvature.
The full one-point measure at \(D>M\) contains a finite mass fixed at
\(\pm1\), while the atom-stripped sector displays Jacobi repulsion and its
finite-rank ramp.
The exact modes record a dimension-theorem intersection of the accessible row
space with the two polarizations of \(J\); the ramp probes correlations in the
remaining orientation.
Exact geometry and chaotic geometry are therefore compatible within
orthogonal sectors of the same degenerate manifold.

The algebraic atom count also strengthens the analysis.  Counts are fixed from
\((D,M)\) before eigenvalues are examined, their numerical positions are
audited afterward, and both full and continuous-sector normalizations are
reported.
